Un long fil rectiligne est parcouru par un courant de \qty{\v{I}}{A}. Un proton ($m_p = \qty{1.673e-27}{kg}$, $e = \qty{1.602e-19}{C}$) se déplace parallèlement au fil, dans le même sens que le courant, à \qty{\v{r}}{cm} du fil et à une vitesse de \qty{\v{v}e5}{m/s}. ($\mu_0 = 4\pi \times \qty{e-7}{T.m/A}$)

\begin{enumerate}[label=\alph*)]
    \item Déterminez le champ magnétique produit par le fil à la position du proton.
    \item Déterminez la grandeur de la force magnétique sur le proton.
    \item Le proton est-il attiré vers le fil ou repoussé ? Justifiez.
    \item Déterminez la grandeur de l'accélération du proton à cet instant.
\end{enumerate}
